\documentclass[11pt,a4paper]{article}
\newcommand{\gradelabel}{Grade 7}
\newcommand{\leconhdr}{Les conditions Si / Alors / Sinon}
% Préambule commun aux plans de cours de robotique (Grades 7 à 12).
% Chaque plan de cours commence par \documentclass[11pt,a4paper]{article}, définit
% \gradelabel et \leconhdr (en-têtes de page), puis % Préambule commun aux plans de cours de robotique (Grades 7 à 12).
% Chaque plan de cours commence par \documentclass[11pt,a4paper]{article}, définit
% \gradelabel et \leconhdr (en-têtes de page), puis % Préambule commun aux plans de cours de robotique (Grades 7 à 12).
% Chaque plan de cours commence par \documentclass[11pt,a4paper]{article}, définit
% \gradelabel et \leconhdr (en-têtes de page), puis \input{../../preambule_robotique.tex}.
\usepackage[T1]{fontenc}
\usepackage[utf8]{inputenc}
\usepackage{lmodern}
\usepackage[french]{babel}
\usepackage[margin=2cm,headheight=14pt]{geometry}
\usepackage[most]{tcolorbox}
\usepackage{xcolor,tabularx,booktabs,array,enumitem,fancyhdr,amssymb,tikz}
\usetikzlibrary{arrows.meta,positioning}

\pagestyle{fancy}
\fancyhf{}
\lhead{\small Robotique -- \gradelabel{} -- Novembre 2026}
\rhead{\small \leconhdr}
\cfoot{\thepage}
\setlist{itemsep=1pt,topsep=2pt,leftmargin=*}
\renewcommand{\arraystretch}{1.35}
\newcolumntype{Y}{>{\raggedright\arraybackslash}X}
\newcommand{\cb}{$\square$\,}
\newcommand{\fl}{$\rightarrow$}
\newcommand{\lig}{\noindent\rule{\linewidth}{0.3pt}\\[5mm]}

\newtcolorbox{etape}[3]{enhanced,breakable,colback=#1!6,colframe=#1!70!black,
  coltitle=white,fonttitle=\bfseries,title={#2\hfill #3},boxrule=0.8pt,arc=2mm,
  before skip=8pt,after skip=8pt}
\newtcolorbox{astuce}[1][Différenciation et besoins particuliers]{breakable,
  colback=yellow!12,colframe=yellow!50!black,fonttitle=\bfseries\small,
  title={#1},boxrule=0.5pt,arc=1mm,before skip=6pt,after skip=6pt}
\newtcolorbox{code}[1][Code]{breakable,colback=black!4,colframe=black!60,
  fonttitle=\bfseries\small,title={#1},boxrule=0.5pt,arc=1mm,fontupper=\ttfamily\small,
  before skip=6pt,after skip=6pt}
% Blocs de code coloré (écrire le code sans accents : comme dans l'éditeur)
\newtcblisting{arduino}[1][Code Arduino]{listing only,breakable,colback=black!4,
  colframe=black!60,fonttitle=\bfseries\small,title={#1},boxrule=0.5pt,arc=1mm,
  before skip=6pt,after skip=6pt,
  listing options={language=C++,basicstyle=\ttfamily\small,columns=fullflexible,
  keepspaces=true,showstringspaces=false,commentstyle=\color{green!40!black}\itshape,
  keywordstyle=\color{blue!70!black}\bfseries}}
\newtcblisting{python}[1][Code Python]{listing only,breakable,colback=black!4,
  colframe=black!60,fonttitle=\bfseries\small,title={#1},boxrule=0.5pt,arc=1mm,
  before skip=6pt,after skip=6pt,
  listing options={language=Python,basicstyle=\ttfamily\small,columns=fullflexible,
  keepspaces=true,showstringspaces=false,commentstyle=\color{green!40!black}\itshape,
  keywordstyle=\color{blue!70!black}\bfseries}}
\newcommand{\entete}[2]{%
  \begin{tcolorbox}[colback=#1!10,colframe=#1!70!black,boxrule=0.8pt,arc=2mm]
  {\Large\bfseries #2}\hfill Nom : \rule{4cm}{0.3pt}\quad Date : \rule{2.2cm}{0.3pt}
  \end{tcolorbox}}

% Bandeau de titre : #1 titre, #2 semaine
\newcommand{\titreblock}[2]{%
\begin{tcolorbox}[colback=teal!10,colframe=teal!70!black,arc=2mm,boxrule=1pt]
\centering
{\LARGE\bfseries #1}\\[2mm]
{\large Robotique -- \gradelabel{} -- Plan de cours de 2 heures}\\[1mm]
Enseignant : Francis Ki \quad|\quad #2
\end{tcolorbox}}

% Tableau du minutage : 6 contenus (accroche, mini-leçon, guidée, autonome, mise en commun, exit ticket)
\newcommand{\minutage}[6]{%
\begin{tabularx}{\linewidth}{>{\bfseries}p{2.2cm}p{2.6cm}Y p{3cm}}
\toprule
Minutes & Étape & Contenu & Organisation\\
\midrule
0 -- 10 & Accroche & #1 & Classe entière\\
10 -- 30 & Mini-leçon & #2 & Classe entière\\
30 -- 60 & Pratique guidée & #3 & Collectif puis binômes\\
60 -- 95 & Pratique autonome & #4 & Individuel sur poste\\
95 -- 110 & Mise en commun & #5 & Classe entière\\
110 -- 120 & Exit ticket & #6 & Individuel sur poste\\
\bottomrule
\end{tabularx}}

% Étayage progressif : #1 aide 1, #2 aide 2, #3 aide 3
\newcommand{\aides}[3]{%
\textbf{Étayage progressif de l'aide.}
\begin{center}
\begin{tabularx}{\linewidth}{>{\bfseries}p{2.4cm}Y p{4.4cm}}
\toprule
Niveau d'aide & Contenu & Qui intervient\\
\midrule
Aide 1 & #1 & Le binôme seul\\
Aide 2 & #2 & Le binôme avec la carte d'aide\\
Aide 3 & #3 & L'enseignant\\
\bottomrule
\end{tabularx}
\end{center}}

% Les trois groupes de la pratique autonome : #1 tâche A, #2 tâche B, #3 tâche C
\newcommand{\groupes}[3]{%
\textbf{Consigne commune\,:} chaque élève ouvre le fichier de son groupe, le renomme avec son prénom et travaille seul sur son poste. L'enseignant passe d'abord auprès du groupe A pendant 10 min.

\medskip
\begin{tabularx}{\linewidth}{>{\bfseries}p{3cm}Y p{2.2cm}}
\toprule
Groupe & Tâche sur poste & Annexe\\
\midrule
A : en dessous du niveau & #1 & Annexe C\\
B : au niveau & #2 & Annexe D\\
C : au-dessus du niveau & #3 & Annexe E\\
\bottomrule
\end{tabularx}}

% Mise en commun : #1 groupe A, #2 groupe B, #3 groupe C, #4 conclusion
\newcommand{\miseencommun}[4]{%
\begin{etape}{teal}{5. Mise en commun}{95 -- 110 min}
\begin{enumerate}
 \item (2 min) Retour au calme, écrans mis en veille.
 \item (9 min) Trois passages courts de 3 min, avec valorisation :
 \begin{itemize}
  \item \textbf{Groupe A} : #1
  \item \textbf{Groupe B} : #2
  \item \textbf{Groupe C} : #3
 \end{itemize}
 \item (4 min) Conclusion : #4
\end{enumerate}
\end{etape}}

% Clôture : #1 question 1, #2 question 2
\newcommand{\cloture}[2]{%
\begin{etape}{gray}{6. Clôture : exit ticket}{110 -- 120 min}
Formulaire en ligne ou papier (annexe G) :
\begin{enumerate}
 \item #1
 \item #2
 \item Auto-évaluation : \emph{Je maîtrise / Je progresse / J'ai besoin d'aide.}
\end{enumerate}
Lire les réponses avant la séance suivante pour ajuster les groupes de besoin.
\end{etape}}

% Grille d'observation : #1 #2 #3 intitulés de colonnes
\newcommand{\grille}[3]{%
\subsection*{Grille d'observation (enseignant)}
\begin{tabularx}{\linewidth}{Y p{2.4cm} p{2.4cm} p{2.4cm} Y}
\toprule
Élève & #1 & #2 & #3 & Groupe suivant\\
\midrule
 & & & & \\ & & & & \\ & & & & \\ & & & & \\ & & & & \\
\bottomrule
\end{tabularx}}

% Signature de l'exit ticket
\newcommand{\autoeval}{Je me sens : \cb Je maîtrise \quad \cb Je progresse \quad \cb J'ai besoin d'aide}
.
\usepackage[T1]{fontenc}
\usepackage[utf8]{inputenc}
\usepackage{lmodern}
\usepackage[french]{babel}
\usepackage[margin=2cm,headheight=14pt]{geometry}
\usepackage[most]{tcolorbox}
\usepackage{xcolor,tabularx,booktabs,array,enumitem,fancyhdr,amssymb,tikz}
\usetikzlibrary{arrows.meta,positioning}

\pagestyle{fancy}
\fancyhf{}
\lhead{\small Robotique -- \gradelabel{} -- Novembre 2026}
\rhead{\small \leconhdr}
\cfoot{\thepage}
\setlist{itemsep=1pt,topsep=2pt,leftmargin=*}
\renewcommand{\arraystretch}{1.35}
\newcolumntype{Y}{>{\raggedright\arraybackslash}X}
\newcommand{\cb}{$\square$\,}
\newcommand{\fl}{$\rightarrow$}
\newcommand{\lig}{\noindent\rule{\linewidth}{0.3pt}\\[5mm]}

\newtcolorbox{etape}[3]{enhanced,breakable,colback=#1!6,colframe=#1!70!black,
  coltitle=white,fonttitle=\bfseries,title={#2\hfill #3},boxrule=0.8pt,arc=2mm,
  before skip=8pt,after skip=8pt}
\newtcolorbox{astuce}[1][Différenciation et besoins particuliers]{breakable,
  colback=yellow!12,colframe=yellow!50!black,fonttitle=\bfseries\small,
  title={#1},boxrule=0.5pt,arc=1mm,before skip=6pt,after skip=6pt}
\newtcolorbox{code}[1][Code]{breakable,colback=black!4,colframe=black!60,
  fonttitle=\bfseries\small,title={#1},boxrule=0.5pt,arc=1mm,fontupper=\ttfamily\small,
  before skip=6pt,after skip=6pt}
% Blocs de code coloré (écrire le code sans accents : comme dans l'éditeur)
\newtcblisting{arduino}[1][Code Arduino]{listing only,breakable,colback=black!4,
  colframe=black!60,fonttitle=\bfseries\small,title={#1},boxrule=0.5pt,arc=1mm,
  before skip=6pt,after skip=6pt,
  listing options={language=C++,basicstyle=\ttfamily\small,columns=fullflexible,
  keepspaces=true,showstringspaces=false,commentstyle=\color{green!40!black}\itshape,
  keywordstyle=\color{blue!70!black}\bfseries}}
\newtcblisting{python}[1][Code Python]{listing only,breakable,colback=black!4,
  colframe=black!60,fonttitle=\bfseries\small,title={#1},boxrule=0.5pt,arc=1mm,
  before skip=6pt,after skip=6pt,
  listing options={language=Python,basicstyle=\ttfamily\small,columns=fullflexible,
  keepspaces=true,showstringspaces=false,commentstyle=\color{green!40!black}\itshape,
  keywordstyle=\color{blue!70!black}\bfseries}}
\newcommand{\entete}[2]{%
  \begin{tcolorbox}[colback=#1!10,colframe=#1!70!black,boxrule=0.8pt,arc=2mm]
  {\Large\bfseries #2}\hfill Nom : \rule{4cm}{0.3pt}\quad Date : \rule{2.2cm}{0.3pt}
  \end{tcolorbox}}

% Bandeau de titre : #1 titre, #2 semaine
\newcommand{\titreblock}[2]{%
\begin{tcolorbox}[colback=teal!10,colframe=teal!70!black,arc=2mm,boxrule=1pt]
\centering
{\LARGE\bfseries #1}\\[2mm]
{\large Robotique -- \gradelabel{} -- Plan de cours de 2 heures}\\[1mm]
Enseignant : Francis Ki \quad|\quad #2
\end{tcolorbox}}

% Tableau du minutage : 6 contenus (accroche, mini-leçon, guidée, autonome, mise en commun, exit ticket)
\newcommand{\minutage}[6]{%
\begin{tabularx}{\linewidth}{>{\bfseries}p{2.2cm}p{2.6cm}Y p{3cm}}
\toprule
Minutes & Étape & Contenu & Organisation\\
\midrule
0 -- 10 & Accroche & #1 & Classe entière\\
10 -- 30 & Mini-leçon & #2 & Classe entière\\
30 -- 60 & Pratique guidée & #3 & Collectif puis binômes\\
60 -- 95 & Pratique autonome & #4 & Individuel sur poste\\
95 -- 110 & Mise en commun & #5 & Classe entière\\
110 -- 120 & Exit ticket & #6 & Individuel sur poste\\
\bottomrule
\end{tabularx}}

% Étayage progressif : #1 aide 1, #2 aide 2, #3 aide 3
\newcommand{\aides}[3]{%
\textbf{Étayage progressif de l'aide.}
\begin{center}
\begin{tabularx}{\linewidth}{>{\bfseries}p{2.4cm}Y p{4.4cm}}
\toprule
Niveau d'aide & Contenu & Qui intervient\\
\midrule
Aide 1 & #1 & Le binôme seul\\
Aide 2 & #2 & Le binôme avec la carte d'aide\\
Aide 3 & #3 & L'enseignant\\
\bottomrule
\end{tabularx}
\end{center}}

% Les trois groupes de la pratique autonome : #1 tâche A, #2 tâche B, #3 tâche C
\newcommand{\groupes}[3]{%
\textbf{Consigne commune\,:} chaque élève ouvre le fichier de son groupe, le renomme avec son prénom et travaille seul sur son poste. L'enseignant passe d'abord auprès du groupe A pendant 10 min.

\medskip
\begin{tabularx}{\linewidth}{>{\bfseries}p{3cm}Y p{2.2cm}}
\toprule
Groupe & Tâche sur poste & Annexe\\
\midrule
A : en dessous du niveau & #1 & Annexe C\\
B : au niveau & #2 & Annexe D\\
C : au-dessus du niveau & #3 & Annexe E\\
\bottomrule
\end{tabularx}}

% Mise en commun : #1 groupe A, #2 groupe B, #3 groupe C, #4 conclusion
\newcommand{\miseencommun}[4]{%
\begin{etape}{teal}{5. Mise en commun}{95 -- 110 min}
\begin{enumerate}
 \item (2 min) Retour au calme, écrans mis en veille.
 \item (9 min) Trois passages courts de 3 min, avec valorisation :
 \begin{itemize}
  \item \textbf{Groupe A} : #1
  \item \textbf{Groupe B} : #2
  \item \textbf{Groupe C} : #3
 \end{itemize}
 \item (4 min) Conclusion : #4
\end{enumerate}
\end{etape}}

% Clôture : #1 question 1, #2 question 2
\newcommand{\cloture}[2]{%
\begin{etape}{gray}{6. Clôture : exit ticket}{110 -- 120 min}
Formulaire en ligne ou papier (annexe G) :
\begin{enumerate}
 \item #1
 \item #2
 \item Auto-évaluation : \emph{Je maîtrise / Je progresse / J'ai besoin d'aide.}
\end{enumerate}
Lire les réponses avant la séance suivante pour ajuster les groupes de besoin.
\end{etape}}

% Grille d'observation : #1 #2 #3 intitulés de colonnes
\newcommand{\grille}[3]{%
\subsection*{Grille d'observation (enseignant)}
\begin{tabularx}{\linewidth}{Y p{2.4cm} p{2.4cm} p{2.4cm} Y}
\toprule
Élève & #1 & #2 & #3 & Groupe suivant\\
\midrule
 & & & & \\ & & & & \\ & & & & \\ & & & & \\ & & & & \\
\bottomrule
\end{tabularx}}

% Signature de l'exit ticket
\newcommand{\autoeval}{Je me sens : \cb Je maîtrise \quad \cb Je progresse \quad \cb J'ai besoin d'aide}
.
\usepackage[T1]{fontenc}
\usepackage[utf8]{inputenc}
\usepackage{lmodern}
\usepackage[french]{babel}
\usepackage[margin=2cm,headheight=14pt]{geometry}
\usepackage[most]{tcolorbox}
\usepackage{xcolor,tabularx,booktabs,array,enumitem,fancyhdr,amssymb,tikz}
\usetikzlibrary{arrows.meta,positioning}

\pagestyle{fancy}
\fancyhf{}
\lhead{\small Robotique -- \gradelabel{} -- Novembre 2026}
\rhead{\small \leconhdr}
\cfoot{\thepage}
\setlist{itemsep=1pt,topsep=2pt,leftmargin=*}
\renewcommand{\arraystretch}{1.35}
\newcolumntype{Y}{>{\raggedright\arraybackslash}X}
\newcommand{\cb}{$\square$\,}
\newcommand{\fl}{$\rightarrow$}
\newcommand{\lig}{\noindent\rule{\linewidth}{0.3pt}\\[5mm]}

\newtcolorbox{etape}[3]{enhanced,breakable,colback=#1!6,colframe=#1!70!black,
  coltitle=white,fonttitle=\bfseries,title={#2\hfill #3},boxrule=0.8pt,arc=2mm,
  before skip=8pt,after skip=8pt}
\newtcolorbox{astuce}[1][Différenciation et besoins particuliers]{breakable,
  colback=yellow!12,colframe=yellow!50!black,fonttitle=\bfseries\small,
  title={#1},boxrule=0.5pt,arc=1mm,before skip=6pt,after skip=6pt}
\newtcolorbox{code}[1][Code]{breakable,colback=black!4,colframe=black!60,
  fonttitle=\bfseries\small,title={#1},boxrule=0.5pt,arc=1mm,fontupper=\ttfamily\small,
  before skip=6pt,after skip=6pt}
% Blocs de code coloré (écrire le code sans accents : comme dans l'éditeur)
\newtcblisting{arduino}[1][Code Arduino]{listing only,breakable,colback=black!4,
  colframe=black!60,fonttitle=\bfseries\small,title={#1},boxrule=0.5pt,arc=1mm,
  before skip=6pt,after skip=6pt,
  listing options={language=C++,basicstyle=\ttfamily\small,columns=fullflexible,
  keepspaces=true,showstringspaces=false,commentstyle=\color{green!40!black}\itshape,
  keywordstyle=\color{blue!70!black}\bfseries}}
\newtcblisting{python}[1][Code Python]{listing only,breakable,colback=black!4,
  colframe=black!60,fonttitle=\bfseries\small,title={#1},boxrule=0.5pt,arc=1mm,
  before skip=6pt,after skip=6pt,
  listing options={language=Python,basicstyle=\ttfamily\small,columns=fullflexible,
  keepspaces=true,showstringspaces=false,commentstyle=\color{green!40!black}\itshape,
  keywordstyle=\color{blue!70!black}\bfseries}}
\newcommand{\entete}[2]{%
  \begin{tcolorbox}[colback=#1!10,colframe=#1!70!black,boxrule=0.8pt,arc=2mm]
  {\Large\bfseries #2}\hfill Nom : \rule{4cm}{0.3pt}\quad Date : \rule{2.2cm}{0.3pt}
  \end{tcolorbox}}

% Bandeau de titre : #1 titre, #2 semaine
\newcommand{\titreblock}[2]{%
\begin{tcolorbox}[colback=teal!10,colframe=teal!70!black,arc=2mm,boxrule=1pt]
\centering
{\LARGE\bfseries #1}\\[2mm]
{\large Robotique -- \gradelabel{} -- Plan de cours de 2 heures}\\[1mm]
Enseignant : Francis Ki \quad|\quad #2
\end{tcolorbox}}

% Tableau du minutage : 6 contenus (accroche, mini-leçon, guidée, autonome, mise en commun, exit ticket)
\newcommand{\minutage}[6]{%
\begin{tabularx}{\linewidth}{>{\bfseries}p{2.2cm}p{2.6cm}Y p{3cm}}
\toprule
Minutes & Étape & Contenu & Organisation\\
\midrule
0 -- 10 & Accroche & #1 & Classe entière\\
10 -- 30 & Mini-leçon & #2 & Classe entière\\
30 -- 60 & Pratique guidée & #3 & Collectif puis binômes\\
60 -- 95 & Pratique autonome & #4 & Individuel sur poste\\
95 -- 110 & Mise en commun & #5 & Classe entière\\
110 -- 120 & Exit ticket & #6 & Individuel sur poste\\
\bottomrule
\end{tabularx}}

% Étayage progressif : #1 aide 1, #2 aide 2, #3 aide 3
\newcommand{\aides}[3]{%
\textbf{Étayage progressif de l'aide.}
\begin{center}
\begin{tabularx}{\linewidth}{>{\bfseries}p{2.4cm}Y p{4.4cm}}
\toprule
Niveau d'aide & Contenu & Qui intervient\\
\midrule
Aide 1 & #1 & Le binôme seul\\
Aide 2 & #2 & Le binôme avec la carte d'aide\\
Aide 3 & #3 & L'enseignant\\
\bottomrule
\end{tabularx}
\end{center}}

% Les trois groupes de la pratique autonome : #1 tâche A, #2 tâche B, #3 tâche C
\newcommand{\groupes}[3]{%
\textbf{Consigne commune\,:} chaque élève ouvre le fichier de son groupe, le renomme avec son prénom et travaille seul sur son poste. L'enseignant passe d'abord auprès du groupe A pendant 10 min.

\medskip
\begin{tabularx}{\linewidth}{>{\bfseries}p{3cm}Y p{2.2cm}}
\toprule
Groupe & Tâche sur poste & Annexe\\
\midrule
A : en dessous du niveau & #1 & Annexe C\\
B : au niveau & #2 & Annexe D\\
C : au-dessus du niveau & #3 & Annexe E\\
\bottomrule
\end{tabularx}}

% Mise en commun : #1 groupe A, #2 groupe B, #3 groupe C, #4 conclusion
\newcommand{\miseencommun}[4]{%
\begin{etape}{teal}{5. Mise en commun}{95 -- 110 min}
\begin{enumerate}
 \item (2 min) Retour au calme, écrans mis en veille.
 \item (9 min) Trois passages courts de 3 min, avec valorisation :
 \begin{itemize}
  \item \textbf{Groupe A} : #1
  \item \textbf{Groupe B} : #2
  \item \textbf{Groupe C} : #3
 \end{itemize}
 \item (4 min) Conclusion : #4
\end{enumerate}
\end{etape}}

% Clôture : #1 question 1, #2 question 2
\newcommand{\cloture}[2]{%
\begin{etape}{gray}{6. Clôture : exit ticket}{110 -- 120 min}
Formulaire en ligne ou papier (annexe G) :
\begin{enumerate}
 \item #1
 \item #2
 \item Auto-évaluation : \emph{Je maîtrise / Je progresse / J'ai besoin d'aide.}
\end{enumerate}
Lire les réponses avant la séance suivante pour ajuster les groupes de besoin.
\end{etape}}

% Grille d'observation : #1 #2 #3 intitulés de colonnes
\newcommand{\grille}[3]{%
\subsection*{Grille d'observation (enseignant)}
\begin{tabularx}{\linewidth}{Y p{2.4cm} p{2.4cm} p{2.4cm} Y}
\toprule
Élève & #1 & #2 & #3 & Groupe suivant\\
\midrule
 & & & & \\ & & & & \\ & & & & \\ & & & & \\ & & & & \\
\bottomrule
\end{tabularx}}

% Signature de l'exit ticket
\newcommand{\autoeval}{Je me sens : \cb Je maîtrise \quad \cb Je progresse \quad \cb J'ai besoin d'aide}


\begin{document}
\titreblock{Les conditions Si / Alors / Sinon}{Semaine du 2 au 6 novembre 2026}

%==================================================================
\section{Informations générales}
\begin{tabularx}{\linewidth}{>{\bfseries}p{3.6cm}Y}
\toprule
Titre & Les conditions Si / Alors / Sinon (Scratch)\\
Classe / effectif & Grade 7, 11 à 20 élèves, niveaux hétérogènes, élèves à besoins particuliers\\
Durée & 2 heures (120 min)\\
Lieu & Salle informatique, un poste par élève\\
Matériel et logiciels &
\begin{itemize}
 \item Un poste par élève avec Scratch 3 (éditeur en ligne scratch.mit.edu ou version hors ligne).
 \item Vidéoprojecteur ou tableau interactif.
 \item Des projets de départ «~à compléter~» préparés par l'enseignant (un fichier \texttt{.sb3} par niveau), à placer dans un dossier partagé.
 \item Facultatif : un compte Scratch par élève pour enregistrer en ligne. Sinon, enregistrer en \texttt{.sb3} sur le poste.
\end{itemize}\\
Pré-requis &
Blocs de mouvement et d'apparence, création d'une animation simple, notion de variable (septembre). Le bloc \texttt{si\ldots{} alors} a été découvert du 19 au 23 octobre.\\
\bottomrule
\end{tabularx}

\begin{astuce}[Préparation avant le cours (environ 40 min)]
\begin{enumerate}
 \item Préparer trois projets de départ : A (chat qui change de couleur, avec blocs à compléter), B (jeu «~plus grand ou plus petit~» sans la condition), C (jeu de devinette vide).
 \item Tester les projets dans la version utilisée par les élèves.
 \item Créer le formulaire de sortie (annexe G) dans Google Forms ou imprimer l'annexe.
 \item Constituer les groupes A, B, C à partir du Map Test d'octobre et des observations.
\end{enumerate}
\end{astuce}

%==================================================================
\section{Objectifs d'apprentissage}
\begin{tcolorbox}[colback=gray!8,colframe=gray!60!black,arc=1mm,boxrule=0.5pt]
\textbf{Objectifs du programme}
\begin{itemize}
 \item Utiliser plusieurs choix dans un programme.
 \item Développer la logique algorithmique.
\end{itemize}
\textbf{À la fin de la séance, l'élève est capable\,:}
\begin{itemize}
 \item d'écrire une condition avec un test (égal, plus grand, plus petit)\,;
 \item d'utiliser les blocs \texttt{si\ldots{} alors} et \texttt{si\ldots{} alors\ldots{} sinon} pour obtenir deux comportements différents\,;
 \item de prédire le résultat d'un programme avant de l'exécuter et d'expliquer l'écart si besoin.
\end{itemize}
\textbf{Compétences numériques\,:} ouvrir et enregistrer un projet, tester et corriger, nommer ses fichiers.
\end{tcolorbox}

%==================================================================
\section{Déroulé minuté (120 min)}
\minutage{Le jeu «~Si\ldots{} alors\ldots~» sans écran}{Condition, test et sinon}{Un jeu «~plus grand ou plus petit~» collectif, puis un deuxième en binômes}{Trois niveaux de besoin}{Valoriser chaque groupe}{Trois questions sur poste}

%------------------------------------------------------------------
\begin{etape}{orange}{1. Accroche : le jeu «~Si\ldots{} alors\ldots~»}{0 -- 10 min}
\textbf{But\,:} faire sentir qu'un programme choisit selon une situation.\\
\textbf{Écrans éteints.}
\begin{enumerate}
 \item (5 min) L'enseignant joue le robot. Il annonce des règles : \emph{Si je lève la main, alors vous vous levez. Sinon vous restez assis.} Faire jouer trois règles, dont une avec deux conditions (\emph{si la carte est rouge}).
 \item (5 min) Question-problème : \textbf{«~Comment un programme sait-il quoi faire quand il y a deux possibilités ?~»} Recueillir des exemples de la vie courante (feu rouge ou vert, porte automatique).
\end{enumerate}
\end{etape}

%------------------------------------------------------------------
\begin{etape}{blue}{2. Mini-leçon : condition, test, sinon}{10 -- 30 min}
\textbf{Support\,:} schéma construit au tableau (annexe A). Quatre temps de 5 min.

\textbf{Temps 1 : une condition.} Une condition est une question dont la réponse est \emph{vrai} ou \emph{faux}. Exemples : \texttt{touche la couleur rouge ?}, \texttt{score > 5}.

\textbf{Temps 2 : le bloc si\ldots{} alors.} Si la condition est vraie, on exécute les blocs à l'intérieur. Sinon, on les saute.

\textbf{Temps 3 : le bloc sinon.} Deux chemins, un seul est exécuté.

\begin{center}
\begin{tikzpicture}[box/.style={draw,rounded corners,thick,minimum height=1cm,text width=3cm,align=center,font=\small},>={Stealth[length=3mm]}]
\node[box,fill=blue!15] (c) {Condition vraie ou fausse ?};
\node[box,fill=green!15,below left=0.8cm and 0.5cm of c] (v) {Vrai : \\ blocs du \texttt{alors}};
\node[box,fill=red!10,below right=0.8cm and 0.5cm of c] (f) {Faux : \\ blocs du \texttt{sinon}};
\draw[->,thick] (c)--(v); \draw[->,thick] (c)--(f);
\end{tikzpicture}
\end{center}

\textbf{Temps 4 : prédire avant d'exécuter.} Écrire au tableau trois petits programmes. Pour chacun, la classe lève 1, 2 ou 3 doigts pour son résultat prévu, puis on teste.

\begin{astuce}
Donner une couleur à chaque bloc (condition en bleu, alors en vert, sinon en rouge) sur tous les supports. Lire chaque programme à voix haute en disant «~si~», «~alors~», «~sinon~».
\end{astuce}
\end{etape}

%------------------------------------------------------------------
\begin{etape}{green}{3. Pratique guidée}{30 -- 60 min}
\textbf{Fiche\,:} annexe B.

\textbf{3a. Exemple collectif sur poste (12 min).} Ouvrir le projet de départ B. Objectif : le chat dit «~Gagné~» si le nombre secret est égal au nombre tapé, «~Perdu~» sinon.
\begin{enumerate}
 \item Lire le programme existant et prédire ce qu'il fait.
 \item Ajouter ensemble le bloc \texttt{si\ldots{} alors\ldots{} sinon} avec le test \texttt{réponse = nombre secret}.
 \item Tester avec une bonne et une mauvaise réponse.
\end{enumerate}

\textbf{3b. Second exemple en binômes (18 min).} Programmer un chat qui dit «~plus petit~» ou «~plus grand~» selon le nombre tapé. Alternance des rôles à mi-temps.

\aides{Projet B corrigé du 3a affiché au tableau et schéma de l'annexe A}{Questions guides : \emph{Quelle est la condition ? Que se passe-t-il si elle est vraie ? Et si elle est fausse ?}}{L'enseignant place le bloc \texttt{si} dans le programme, le binôme complète}

\begin{astuce}
Proposer un projet où les blocs sont déjà présents mais mélangés pour les élèves qui ont peur d'un écran vide. Laisser le temps de tester plusieurs fois.
\end{astuce}
\end{etape}

%------------------------------------------------------------------
\begin{etape}{violet}{4. Pratique autonome différenciée}{60 -- 95 min}
\groupes{Projet à compléter : les blocs sont donnés, il faut les ranger dans l'ordre dans un \texttt{si\ldots{} alors\ldots{} sinon}, avec des étiquettes de couleur (chat qui change de couleur selon la touche pressée).}{Programmer seul un jeu de devinette : un nombre secret, trois essais, des messages «~plus grand~», «~plus petit~», «~gagné~» avec des conditions imbriquées.}{Concevoir un jeu avec deux conditions combinées (\texttt{et}, \texttt{ou}) : par exemple un personnage qui évite un obstacle seulement si la touche flèche est pressée \emph{et} la couleur est rouge. Documenter le programme avec des commentaires.}

\medskip
\textbf{Plan B si un élève termine tôt\,:} corriger le projet d'un camarade dont l'erreur a été préparée par l'enseignant.
\end{etape}

%------------------------------------------------------------------
\miseencommun{un élève lit à voix haute son programme avec «~si~», «~alors~», «~sinon~». Valoriser la lecture dans le bon ordre.}{un binôme projette son jeu et explique ses conditions. Valoriser la clarté.}{un élève présente sa condition combinée. Valoriser la complexité maîtrisée.}{\textit{une condition donne deux chemins, un seul est exécuté.} Annoncer que la leçon suivante utilisera des boucles.}

\cloture{Quel bloc faut-il utiliser pour obtenir deux comportements différents selon une condition ?}{Que fait ce programme : si 3 > 5 alors dire «~A~» sinon dire «~B~» ?}

%==================================================================
% ANNEXES
%==================================================================
\newpage
\section*{Annexe A : Support résumé de la mini-leçon}
\entete{blue}{Les conditions Si / Alors / Sinon}

\begin{tcolorbox}[colback=yellow!10,colframe=yellow!50!black]
\textbf{Condition} : une question dont la réponse est \emph{vrai} ou \emph{faux}.\\
\textbf{Si\ldots{} alors} : exécute des blocs seulement si la condition est vraie.\\
\textbf{Si\ldots{} alors\ldots{} sinon} : choisit entre deux chemins.
\end{tcolorbox}

\begin{tabularx}{\linewidth}{>{\bfseries}p{3.4cm}Y Y}
\toprule
Test & Exemple Scratch & Vrai quand\ldots\\
\midrule
Égal & \texttt{réponse = 7} & la réponse vaut 7\\
Plus grand & \texttt{score > 5} & le score dépasse 5\\
Plus petit & \texttt{vie < 3} & la vie est inférieure à 3\\
Touche & \texttt{touche espace pressée ?} & on appuie sur espace\\
\bottomrule
\end{tabularx}

\begin{code}[Modèle]
si <condition> alors\\
\hspace*{1cm}blocs si vrai\\
sinon\\
\hspace*{1cm}blocs si faux
\end{code}

%==================================================================
\newpage
\section*{Annexe B : Fiche de pratique guidée}
\entete{green}{Pratique guidée : si / alors / sinon}

\subsection*{3a. Exemple collectif : gagné ou perdu}
Condition choisie : \dotfill\\
Si vrai, le chat dit : \dotfill\\
Si faux, le chat dit : \dotfill

\textbf{Test 1} (bonne réponse) : résultat \dotfill\quad \textbf{Test 2} (mauvaise réponse) : résultat \dotfill

\subsection*{3b. En binôme : plus grand ou plus petit}
\begin{tabularx}{\linewidth}{>{\bfseries}p{4cm}Y Y}
\toprule
Étape & Ce que je fais & Mon test\\
\midrule
1. Nombre secret & & \\
2. Demander un nombre & & \\
3. Condition & & \\
4. Messages & & \\
\bottomrule
\end{tabularx}

\begin{tcolorbox}[colback=yellow!10,colframe=yellow!50!black,title={Cartes d'aide à utiliser dans l'ordre},fonttitle=\bfseries\small]
\textbf{Aide 1 :} regarde le programme du 3a.\\
\textbf{Aide 2 :} Quelle est la condition ? Que se passe-t-il si elle est vraie ? Et si elle est fausse ?\\
\textbf{Aide 3 :} demande à l'enseignant de placer le bloc \texttt{si} avec toi.
\end{tcolorbox}

%==================================================================
\newpage
\section*{Annexe C : Niveau A (en dessous du niveau)}
\entete{orange}{Niveau A : je range les blocs}

\textbf{Projet de départ :} «~Chat qui change de couleur~» (blocs donnés, à ranger).

\subsection*{Activité 1 : relie chaque situation à sa condition}
\begin{tabularx}{\linewidth}{Y Y}
\toprule
\textbf{Situation} & \textbf{Condition à utiliser}\\
\midrule
Le chat change de couleur si on appuie sur espace & \texttt{touche espace pressée ?}\\
Le chat dit «~Gagné~» si le score est plus grand que 10 & \texttt{score > 10}\\
Le chat avance si la réponse est 5 & \texttt{réponse = 5}\\
\bottomrule
\end{tabularx}

\subsection*{Activité 2 : mets les blocs dans l'ordre (numérote)}
\dotfill{} \texttt{sinon : dire «~Perdu~»}\\
\dotfill{} \texttt{si réponse = 7 alors : dire «~Gagné~»}\\
\dotfill{} \texttt{demander «~Quel nombre ?~»}

\subsection*{Activité 3 : prédis puis teste}
Que se passe-t-il si la réponse est 7 ? \dotfill\quad Si elle est 4 ? \dotfill

%==================================================================
\newpage
\section*{Annexe D : Niveau B (au niveau)}
\entete{blue}{Niveau B : je programme un jeu de devinette}

\textbf{Cahier des charges :} le programme choisit un nombre secret entre 1 et 10, laisse trois essais et répond «~plus grand~», «~plus petit~» ou «~gagné~».

\begin{tabularx}{\linewidth}{>{\bfseries}p{4cm}Y}
\toprule
Étape & Mon programme (blocs principaux)\\
\midrule
Nombre secret & \\
Essai 1 & \\
Essai 2 & \\
Essai 3 & \\
\bottomrule
\end{tabularx}

\subsection*{Questions}
\begin{enumerate}
 \item Combien de conditions as-tu utilisées ? \dotfill
 \item Quel cas n'est pas géré dans ton programme ? \dotfill
\end{enumerate}

%==================================================================
\newpage
\section*{Annexe E : Niveau C (au-dessus du niveau)}
\entete{violet}{Niveau C : conditions combinées}

\textbf{Mission :} crée un jeu où un personnage doit éviter un obstacle.

\begin{enumerate}
 \item Le personnage saute seulement si la touche flèche est pressée \textbf{et} qu'il touche le sol.
 \item Si l'obstacle touche le personnage, alors le score diminue, sinon le score augmente.
 \item Ajoute des commentaires (blocs de commentaire) qui expliquent chaque condition.
\end{enumerate}

\begin{tabularx}{\linewidth}{>{\bfseries}p{4cm}Y}
\toprule
Condition & Rôle dans le jeu\\
\midrule
 & \\ & \\ & \\
\bottomrule
\end{tabularx}

\textbf{Défi bonus :} ajoute un niveau de difficulté : l'obstacle accélère si le score dépasse 10.

%==================================================================
\newpage
\section*{Annexe F : Corrections et attendus (enseignant)}

\subsection*{Annexe B}
3a : condition \texttt{réponse = nombre secret} ; vrai : «~Gagné~» ; faux : «~Perdu~». 3b : \texttt{si réponse = nombre secret alors dire «~gagné~» sinon} (puis, \emph{à l'intérieur du sinon}) \texttt{si réponse < nombre secret alors dire «~plus grand~» sinon dire «~plus petit~»}. Scratch n'a pas de bloc «~sinon si~» : on imbrique un deuxième \texttt{si}.

\subsection*{Annexe C}
Activité 1 : relier dans l'ordre du tableau. Activité 2 : 1 demander ; 2 si réponse = 7 ; 3 sinon. Activité 3 : «~Gagné~» ; «~Perdu~».

\subsection*{Annexe D}
Programme attendu : nombre secret aléatoire, boucle de trois essais (ou trois blocs), trois cas gérés (plus petit, plus grand, égal). Cas oublié courant : réponse non numérique.

\subsection*{Annexe E}
Conditions attendues : \texttt{touche flèche pressée} \textbf{et} \texttt{touche le sol}, \texttt{touche l'obstacle ?} avec sinon.\\
\textbf{Critères de réussite :} deux conditions correctes, un comportement différent selon le résultat, des commentaires utiles.

%==================================================================
\newpage
\section*{Annexe G : Exit ticket et grille d'observation}
\entete{gray}{Exit ticket : si / alors / sinon}

\begin{enumerate}
 \item Quel bloc faut-il utiliser pour obtenir deux comportements différents selon une condition ?\\ \lig
 \item Que fait ce programme : si 3 > 5 alors dire «~A~» sinon dire «~B~» ?\\ \lig
 \item \autoeval
\end{enumerate}
\textbf{Attendus :} 1. \texttt{si\ldots{} alors\ldots{} sinon}. 2. Il dit «~B~» car 3 n'est pas plus grand que 5.

\bigskip
\grille{Condition écrite}{Sinon utilisé}{Test réalisé}

\end{document}
